% Concept / architecture figure: paradigm inversion + four-stage pipeline.
% Self-contained TikZ (no external assets).
\begin{figure}[t]
\centering
\resizebox{\columnwidth}{!}{%
\begin{tikzpicture}[
  font=\small,
  >={Stealth[round]},
  box/.style={rounded corners=2pt, draw, thick, minimum height=8mm,
              minimum width=15mm, align=center, inner sep=3pt},
  inbox/.style={box, fill=blue!10, draw=blue!55},
  outbox/.style={box, fill=red!10, draw=red!55},
  stage/.style={box, fill=green!12, draw=green!50!black, minimum width=17mm},
  lbl/.style={font=\footnotesize\itshape, text=black!60},
  ar/.style={->, thick},
]

% ---- top: classical (input-side) ----
\node[inbox] (cin) {Input\\covering\\array};
\node[box, fill=black!6, right=8mm of cin] (csut) {SUT\\$f:X\!\to\!\mathcal{P}(Y)$};
\node[outbox, right=8mm of csut] (cout) {Outputs\\(uncontrolled)};
\draw[ar, blue!55] (cin) -- (csut);
\draw[ar, black!45] (csut) -- (cout);
\node[lbl, above=1mm of csut] {classical: cover inputs, observe outputs};

% ---- divider label ----
\node[lbl, below=4mm of csut, text=red!70!black, font=\footnotesize\bfseries]
  (inv) {paradigm inversion};

% ---- bottom: reverse (output-side) four stages ----
\node[stage, below=10mm of cin, xshift=-2mm] (s1)
  {\textbf{1} Feasibility\\$\hat f(x)\!\in\!W$};
\node[stage, right=5mm of s1] (s2)
  {\textbf{2} Output\\cov.\ array\\$\mathrm{OCA}$};
\node[stage, right=5mm of s2] (s3)
  {\textbf{3} Inverse\\map\\$\min_x L(x;w^{*})$};
\node[stage, right=5mm of s3] (s4)
  {\textbf{4} Prioritise\\$\mathrm{OCov}_s$};
\draw[ar, green!45!black] (s1) -- (s2);
\draw[ar, green!45!black] (s2) -- (s3);
\draw[ar, green!45!black] (s3) -- (s4);

% feedback: covered output behaviour -> concrete input (the inversion)
\draw[ar, red!60, dashed] (s3.south) .. controls +(0,-7mm) and +(0,-7mm) ..
  node[lbl, below, text=red!70!black] {each covered output behaviour $\Rightarrow$ a concrete input test} (s1.south);

\end{tikzpicture}%
}
\caption{Classical combinatorial testing (top) covers \emph{input} interactions and
observes whatever outputs result. \rnwise{} (bottom) inverts this: it covers a
$t$-way array over \emph{abstracted output behaviours} and recovers concrete
inputs by gradient-free inverse mapping. The four stages are agnostic to the
system's internals, so the same pipeline applies to ML and quantum systems.}
\label{fig:concept}
\end{figure}
